\cotefile{85}
\scanpage{86}
une \unsure{décomposition} en produit $\frac12$ direct
\[
L=G*\gamma\simeq G\underset{1/2}{\cdot}\unsure{E|\mathcal E}
\]
On va interpréter $E$ en notant qu'on a un hom (\unsure{inj|inj.})
\[
\gamma \overset{\rho}{\hookrightarrow} E
\]
et une op. de $G$ sur $E$, d'où pour $\sigma\in G$ un hom
\[
\gamma\overset{{}^{\sigma}\rho}{\unsure{\mathrel{\relbar}|\longrightarrow}}E\qquad
\deleted{\unsure{{}^{\sigma}\rho(x)={}^{\sigma}\rho(x)\,\rho}}
\inserted{{}^{\sigma}\rho=\sigma_E\circ\rho.}
\]
La famille des hom de groupes permet donc de définir un hom
\[
\underbrace{\gamma^{(*G)}}_{\text{puiss. libre}}\longrightarrow E
\]
\deleted{et une op. de}

D'ailleurs l'op. de $G$ sur l'ens d'indices $G$ permet de faire opérer $G$ sur la puiss. libre $\gamma^{(*G)}$, et \deleted{\unsure{par}} on voit que l'hom précédent est compatible aux op. de $G$. \deleted{En fait, on voit \unsure{trivialemt|tout de suite} que c'est un} \inserted{Donc on trouve}
\[
G\underset{1/2}{\cdot}\bigl(\gamma^{*G}\bigr)\longrightarrow L\simeq
G\underset{1/2}{\cdot}E
\]
On voit de suite que cet hom de groupes est un \uline{iso} (cf propriété universelle) donc
\[
\gamma^{(*G)}\xrightarrow{\sim}E\qquad \text{\uuline{iso}}
\]
% archive sheet number: 85
